\documentclass[10pt, letterpaper]{article}

\usepackage[
    ignoreheadfoot,
    top=0.8cm, bottom=0.8cm, left=1.2cm, right=1.2cm,
    footskip=1.0 cm,
    ]{geometry}
\usepackage{titlesec}
\usepackage{tabularx}
\usepackage{array}
\usepackage[dvipsnames]{xcolor}
\definecolor{primaryColor}{RGB}{0, 0, 0}
\usepackage{enumitem}
\usepackage{fontawesome5}
\usepackage{amsmath}
\usepackage[
    pdftitle={Andreas Tzitzikas's Hardware Engineering Resume},
    pdfauthor={Andreas Tzitzikas},
    colorlinks=true,
    urlcolor=primaryColor
]{hyperref}
\usepackage[pscoord]{eso-pic}
\usepackage{calc}
\usepackage{bookmark}
\usepackage{lastpage}
\usepackage{changepage}
\usepackage{paracol}
\usepackage{ifthen}
\usepackage{needspace}
\usepackage{iftex}
\usepackage{datetime}

\ifPDFTeX
    \input{glyphtounicode}
    \pdfgentounicode=1
    \usepackage[T1]{fontenc}
    \usepackage[utf8]{inputenc}
    \usepackage{lmodern}
\fi

\usepackage{charter}

\raggedright
\AtBeginEnvironment{adjustwidth}{\partopsep0pt}
\pagestyle{empty}
\setcounter{secnumdepth}{0}
\setlength{\parindent}{0pt}
\setlength{\topskip}{0pt}
\setlength{\columnsep}{0.2cm}
\pagenumbering{gobble}

\titleformat{\section}{\needspace{4\baselineskip}\bfseries\large}{}{0pt}{}[\vspace{1pt}\titlerule]

\titlespacing{\section}{
    -1pt
}{
    0.2 cm
}{
    0.1 cm
}

\renewcommand\labelitemi{$\vcenter{\hbox{\small$\bullet$}}$}
\newenvironment{highlights}{
    \begin{itemize}[
        topsep=0.05 cm,
        parsep=0.05 cm,
        partopsep=0pt,
        itemsep=0pt,
        leftmargin=0 cm + 10pt
    ]
}{
    \end{itemize}
}

\newenvironment{onecolentry}{
    \begin{adjustwidth}{
        0 cm + 0.00001 cm
    }{
        0 cm + 0.00001 cm
    }
}{
    \end{adjustwidth}
}

\newenvironment{twocolentry}[2][]{
    \onecolentry
    \def\secondColumn{#2}
    \setcolumnwidth{\fill, 4.5 cm}
    \begin{paracol}{2}
}{
    \switchcolumn \raggedleft \secondColumn
    \end{paracol}
    \endonecolentry
}

\newenvironment{header}{
    \setlength{\topsep}{0pt}\par\kern\topsep\centering\linespread{1.5}
}{
    \par\kern\topsep
}

\let\hrefWithoutArrow\href

\begin{document}
    \newcommand{\AND}{\unskip
        \cleaders\copy\ANDbox\hskip\wd\ANDbox
        \ignorespaces
    }
    \newsavebox\ANDbox
    \sbox\ANDbox{$|$}

    \begin{header}
        \fontsize{15 pt}{15 pt}\selectfont Andreas Tzitzikas

        \vspace{3 pt}

        \normalsize
        \mbox{College Park, MD}
        \kern 3.0 pt
        \AND
        \kern 3.0 pt
        \mbox{U.S. Citizen}
        \kern 3.0 pt
        \AND
        \kern 3.0 pt
        \mbox{\hrefWithoutArrow{mailto:andreas.tz.work@gmail.com}{andreas.tz.work@gmail.com}}
        \kern 3.0 pt
        \AND
        \kern 3.0 pt
        \mbox{\hrefWithoutArrow{https://linkedin.com/in/andreas-tzitzikas}{linkedin.com/in/andreas-tzitzikas}}
        \kern 3.0 pt
        \AND
        \kern 3.0 pt
        \mbox{\hrefWithoutArrow{https://github.com/Wafer0}{github.com/Wafer0}}
    \end{header}

    \vspace{3 pt}

    \section{Education}

    \vspace{0.10cm}
    \begin{twocolentry}{Expected May 2027}
        \textbf{University of Maryland}, College Park, MD
    \end{twocolentry}
    \begin{twocolentry}{\textbf{GPA: 4.00}}
        \textit{Master of Science in Computer Engineering}
    \end{twocolentry}
    \vspace{0.05cm}
    \begin{twocolentry}{Expected May 2026}
        \textit{Bachelor of Science in Computer Engineering} (GPA: 3.88)
    \end{twocolentry}
    \vspace{0.05cm}
    \begin{onecolentry}
        \textbf{Relevant Coursework:} Digital Computer Design (ENEE646), Digital CMOS VLSI Design (ENEE640), Microprocessors (ENEE440), Computer Organization (ENEE350), Digital Circuits Lab (ENEE245), PCB Design (ENEE459A), Operating Systems (ENEE447), Reverse Eng. \& HW Security (ENEE459B)
    \end{onecolentry}

    \section{Projects}

        \begin{twocolentry}{
            Aug 2025 – Dec 2025
        }
            \textbf{Dual-Host S2-LP Radio Shield | KiCad, RF Design, Multi-Platform Firmware}
        \end{twocolentry}
        \vspace{0.05 cm}
        \begin{onecolentry}
            \begin{highlights}
                \item Designed dual-host development board for \textbf{S2-LP} sub-1 GHz transceiver with \textbf{STM32} and \textbf{RP2350} compatibility, implementing complete RF circuit with matching network and crystal oscillator for optimal radio performance.
                \item Designed a novel "\textbf{Shadow Mode}" debugging architecture allowing the \textbf{STM32} to snoop \textbf{SPI} traffic between the \textbf{RP2350} and \textbf{S2-LP} radio, facilitating real-time protocol analysis without bus contention.
                \item Created comprehensive firmware suite for both platforms including SPI drivers, radio configuration, packet transmission, and bus monitoring capabilities with detailed technical documentation.
            \end{highlights}
        \end{onecolentry}

        \vspace{0.10 cm}
        \begin{twocolentry}{
            Jan 2025 – May 2025
        }
            \textbf{STM32G4 Bare-Metal Peripheral Integration | ARM Assembly, Embedded C}
        \end{twocolentry}
        \vspace{0.05 cm}
        \begin{onecolentry}
            \begin{highlights}
                \item Implemented \textbf{bare-metal} firmware with direct register-level programming for \textbf{GPIO}, \textbf{DAC}, \textbf{ADC}, \textbf{TIM2} (PWM/PFM), and \textbf{SPI2} flash memory drivers on \textbf{STM32G4} microcontroller.
                \item Wrote \textbf{bare-metal} drivers (\textbf{ARM Assembly}/C) for a resource-constrained \textbf{STM32G4}; implemented a multiplexed interrupt handler to service \textbf{4 distinct peripherals} via a single hardware IRQ, reducing context-switching overhead.
                \item Developed comprehensive testing framework ensuring real-time performance, peripheral synchronization, and robust error handling in resource-constrained embedded environment.
            \end{highlights}
        \end{onecolentry}

        \vspace{0.10 cm}
        \begin{twocolentry}{
            Nov 2025 - Dec 2025 } \textbf{Tomasulo Out-of-Order RISC-V CPU | RTL Design, FPGA Implementation, Verification}
        \end{twocolentry}
        \vspace{0.05 cm}
        \begin{onecolentry}
            \begin{highlights}
                \item Architected complex RTL system using \textbf{Tomasulo's Algorithm} with register renaming, \textbf{8-entry reservation stations}, and \textbf{16-entry reorder buffer} for advanced hardware design.
                \item Achieved \textbf{36\% higher operating frequency} (\textbf{69.7MHz} vs 51.2MHz) and \textbf{40-150\% IPC improvement} compared to in-order implementation through hardware optimization and performance analysis.
                \item Implemented sophisticated branch prediction, speculative execution, and precise exception handling ensuring robust hardware behavior across complex computational workloads.
                \item Synthesized complex digital system with \textbf{1,652 standard cells} achieving \textbf{zero DRC violations} and comprehensive functional verification for reliable hardware operation.
            \end{highlights}
        \end{onecolentry}

        \vspace{0.10 cm}

        \begin{twocolentry}{
            Fall 2025
        }
            \textbf{Polynomial Evaluation Accelerator | SystemVerilog, Digital Design}
        \end{twocolentry}
        \vspace{0.05 cm}
        \begin{onecolentry}
            \begin{highlights}
                \item Architected dataflow-based accelerator with 5 finite state machines managing FIFO handshaking and resource sharing across 9 processing modules for high-throughput computation.
                \item Designed dual-port coefficient memory with auto-padding logic enabling efficient batch processing of polynomial evaluations with minimal memory access latency.
                \item Achieved 100\% functional coverage through comprehensive testbench development with 15+ complex test cases validated against reference C implementation.
            \end{highlights}
        \end{onecolentry}

    \section{Experience}

        \begin{twocolentry}{
            Jun 2025 – Aug 2025
        }
            \textbf{Embedded Systems Intern}, Alchemity -- College Park, MD
        \end{twocolentry}
        \vspace{0.05 cm}
        \begin{onecolentry}
            \begin{highlights}
                \item Accelerated validation of a hybrid solid oxide fuel cell reactor by developing a full-stack automation suite (\textbf{Rust}, \textbf{Python}) to parallelize testing, converting multi-day manual experiments into overnight runs to significantly increase throughput.
                \item Engineered real-time \textbf{STM32} firmware to precisely control material synthesis within modular reactors, featuring non-blocking motor/relay drivers and \textbf{SPI} peripheral management.
                \item Built end-to-end control interfaces including desktop GUIs, embedded display drivers, and CLI tools to streamline workflows between hardware and software teams.
            \end{highlights}
        \end{onecolentry}

        \vspace{0.10 cm}

    \section{Skills}
        \begin{onecolentry}
            \textbf{Languages \& HDLs:} SystemVerilog, Verilog, C/C++, ARM Assembly, Python \\
            \textbf{Hardware Design:} Digital Logic Design, PCB Design, Microcontroller Programming, FPGA Implementation, VLSI Design \\
            \textbf{Embedded Systems:} STM32 Bare-Metal Firmware, Real-Time Control, SPI, I2C, UART, GPIO, ADC/DAC, PWM \\
            \textbf{EDA Tools:} KiCad, Xilinx Vivado, Icarus Verilog, Verible, GTKWave, LPKF Prototyping, Altium Designer \\
            \textbf{Verification:} Functional Verification, Testbench Development, Static Timing Analysis, Linux
        \end{onecolentry}

\end{document}
